\begin{derivation}[盒态与连续态的模式耦合归一化]
先用单位归一周期盒离散化连续动量谱，再由单个动量格点的体积恢复连续核 $\mathcal G^{(\lambda)}_{\vb{p}'s',\vb{p}s}$。盒体积 $V$ 只承担正则化作用；在 $V\to\infty$ 时，它必须在“态密度 $\times$ 矩阵元”的组合中消去。取 $V=L_xL_yL_z$，允许动量为

\begin{equation*}
\vb{p}_{\vb{n}}
=2\pi\hbar\left(\frac{n_x}{L_x},\frac{n_y}{L_y},\frac{n_z}{L_z}\right),
\end{equation*}
盒态归一化为
\begin{equation*}
{}_V\!\langle\vb{p}_{\vb{n}'},s'|\vb{p}_{\vb{n}},s\rangle_V
=\delta_{\vb{n}'\vb{n}}\delta_{s's}.
\end{equation*}
当 $V\to\infty$ 时，每个动量格点占据体积 $\Delta^3p=(2\pi\hbar)^3/V$，所以
\begin{equation*}
\sum_{\vb{n}}F(\vb{p}_{\vb{n}})
\longrightarrow
\frac{V}{(2\pi\hbar)^3}\int\dd^3 p\,F(\vb{p}).
\end{equation*}
与连续约定
$\langle\vb{p}',s'|\vb{p},s\rangle=\delta_{s's}\delta^{(3)}(\vb{p}'-\vb{p})$
相容的态变换是
\begin{equation*}
|\vb{p}_{\vb{n}},s\rangle_V
=\sqrt{\Delta^3p}\,|\vb{p}_{\vb{n}},s\rangle.
\end{equation*}
因此盒中一量子顶点矩阵元
\begin{equation*}
\mathcal G^{(\lambda,V)}_{\vb{n}'s',\vb{n}s}
=\frac1\hbar{}_V\!\langle\vb{p}_{\vb{n}'},s'|
\delta h_\lambda|\vb{p}_{\vb{n}},s\rangle_V
\end{equation*}
与连续核之间满足
\begin{equation*}
\mathcal G^{(\lambda,V)}_{\vb{n}'s',\vb{n}s}
=\Delta^3p\,
\mathcal G^{(\lambda)}_{\vb{p}_{\vb{n}'}s',\vb{p}_{\vb{n}}s}
\end{equation*}
（若初末态使用不同离散化体积，相应平方根因子分别保留；最终物理组合不变）。态展开给出 $C_{\vb{n}}=\sqrt{\Delta^3p}\,C(\vb{p}_{\vb{n}})$。将离散 Schr\"odinger 方程两侧同除输出态的 $\sqrt{\Delta^3p}$，耦合项成为
\begin{align*}
\frac{1}{\sqrt{\Delta^3p}}\sum_{\vb{n}}
\mathcal G^{(\lambda,V)}_{\vb{n}'\vb{n}}C_{\vb{n}}
&=\sum_{\vb{n}}\Delta^3p\,
\mathcal G^{(\lambda)}_{\vb{p}_{\vb{n}'}\vb{p}_{\vb{n}}}C(\vb{p}_{\vb{n}})\\
&\longrightarrow\int\dd^3p\,
\mathcal G^{(\lambda)}_{\vb{p}'\vb{p}}C(\vb{p}).
\end{align*}
所以盒体积在“态密度 $\times$ 盒矩阵元”这一组合中完全消去。连续核的量纲关系是 $[\mathcal G_{\vb{p}'\vb{p}}\,\dd^3p]={\rm s^{-1}}$，与积分测度组合后才具有频率量纲。连续谱中的 $\delta$ 归一化、态密度与耦合矩阵元因此应一起进入物理概率。
\end{derivation}
